% !TEX root = ../../../main.tex
% Block card: detection.sax_psax -- "Symbolic encoding (pSAX)" (Part I, the blocks chapter).
\cardfield{What it does}
Turns a span into a string of letters, one per piece of time, using a fixed number of
letters whose boundaries are fitted to how the signal's piece-averages are actually spread.
(How the SAX family works is explained on the card for \blockref{detection.sax_csax}.)

\cardfield{When to use it}
When you want a level-based string, as with cSAX, but with a number of letters you choose
and always get. It does not assume the signal is bell-shaped: the letters go where the
averages are.

\cardfield{What goes in, what comes out}
In goes a span of one channel. Out comes a string with one letter per piece, \emph{a} for
the lowest level. \blockref{detection.symbol_search} can then search it.

\cardfield{Settings that matter}
\emph{alphabet\_size} is the number of letters (default 8). The piece length is set as for
cSAX: \emph{segment\_mode} picks the live control, by default 20\,s per letter, and the
block page shows the same live read-outs (letters produced, samples dropped and so on)
before you run.

\cardfield{Watch out for}
Above 26 letters the later letters are spelt with two characters (\emph{aa}, \emph{ab}\ldots),
which breaks \blockref{detection.symbol_search}: keep \emph{alphabet\_size} at 26 or below
if you will search the string. On a span that is mostly noise the letters split the noise
finely and still look meaningful. A missing sample (NaN) anywhere in the span makes the
block fail.
